\paragraph{Theorem D.2.37 (Proved: a uniform weighted algorithm on all natural profiles).}
Choose any of the $172$ natural initial rows of D.2.22,
with $P=aG,Q=bG$, and any integer $d>0$ satisfying
\[
                         \gcd(a,d)=1.                    \tag{D.WP1}
\]
Iterate its actual contracting translation $d$ times.
Thus the new translation is $dbG$ and its linearization
is the norm of the original one under the first $d$
translates. Keep the same unimodular Mumford fans,
quotiented by the new period lattice.

Every original/Mumford choice represented by D.2.33,
every free finite-good additive logarithm, every number
of free regular logarithms and every integer peripheral
clutch is covered by the following exact integer
algorithm. Star caps are retained original. The
$495$ initial presentation patterns provide the finite
input; there is no upper bound on $d$, logarithmic
orders or affine coordinates.

Write $A_i$ for the recorded four-dimensional
monodromies in the basis $(e_1,e_2,\delta,\gamma)$.
Replace them by
\[
 A_i(d)=S_d^{-1}A_iS_d,\qquad
 S_d=\operatorname{diag}(1,1,1,d).
                                                               \tag{D.WP2}
\]
Equivalently multiply their first three gamma-column
entries by $d$; every other entry is unchanged.
The primitive common alternating functional is
\[
 \omega_d=a e_1^*e_2^*-db\,\delta^*\gamma^*.                \tag{D.WP3}
\]
Insert identity matrices for regular logarithms.
An additional ordinary order-one logarithm with power
$-e\gamma$ realizes the clutch integer $e$.

Let $N$ be the number of sites, with
$A_0(d)\cdots A_{N-1}(d)=I$, and compute the integral
kernel
\[
 K=\ker\bigl[(A_0(d)-I)\ \cdots\ (A_{N-2}(d)-I)\bigr].
                                                               \tag{D.WP4}
\]
The target of the degree-two presentation is
$\mathbb Zs\oplus K$. For a fixed vector at site $i$,
let $J_i$ denote its usual oriented peripheral
insertion in $K$: the inclusion in block $i$ for
$i<N-1$, and, at the last site, the blocks
$-A_{j+1}(d)\cdots A_{N-2}(d)$.
All these maps are evaluated in a primitive integral
basis of (D.WP4).

At a finite-good site of linear order $h$, compute
a primitive fixed basis $F:\mathbb Z^2\to\Lambda$
and the norm $N_h=I+\cdots+A^{h-1}$. If $H$ is an
integral basis matrix for the image of
$F^{-1}N_h:\Lambda\to\mathbb Z^2$, the projected lattice
is $F(H/h)\mathbb Z^2$. For any invariant power
$v=Fz$ and order $m=hk$, put
\[
 \begin{gathered}
 w=hH^{-1}z\in\mathbb Z^2,\qquad t=\gcd(m,w_1,w_2),\\
 \text{the action is free iff}\quad
 \gcd(m,v_1,v_2,v_3,v_4)=1,\quad h\mid m/t.
 \end{gathered}
                                                               \tag{D.WP5}
\]
For every free such logarithm include the five columns
\[
 \bigl(-\omega_d(\lambda,v),\,J_iN_m\lambda\bigr)
       \quad(\lambda=e_1,e_2,\delta,\gamma),\qquad
                         (0,J_iv),                       \tag{D.WP6}
\]
where $N_m=(m/h)N_h$. At a regular logarithm use the
same formula with $A=I,N_m=mI$ and the sole freeness
condition $\gcd(m,v_1,v_2,v_3,v_4)=1$.

At every unchanged original or Mumford cap include
the full marked fixed-normal columns
\[
                       (0,J_iF_i),\qquad
                   F_i\text{ a basis of }\ker(A_i(d)-I).
                                                               \tag{D.WP7}
\]
For a nonzero-Chern original also include
$(dh_i,0)$, where $h_i$ is its recorded
$b/q_i^c$ or $b/n_i$ from D.2.32--33.
These are the entire global images of the boundary
kernels, including any torsion in a boundary.
Let $B$ be the resulting integer matrix and
$T=\operatorname{Tor}\operatorname{coker}B$.

For the component correction define $\beta$ as follows.
A zero-free rank-one site with $P$ narrow contributes
$\gcd(d,n_i)-1$; if $P$ is non-narrow, its marked
Mumford cap contributes zero. A zero-degree original
additive cap contributes $\rho_i(d)-1$, where
$\rho_i(d)$ is the number of component orbits under
$Q^d$, computed by raising the recorded permutation
to the $d$th power. A nonzero-Chern original contributes
its unchanged $s_i-2$ or $n_i-2$.
Let $L$ count these last originals. The weighted
rank-two zero has $d$ normalized components and is
not included in $\beta$.

The full fundamental group is the Seifert presentation
of D.2.23 with the actual gamma coordinates of the
power vectors, the displayed orders and clutch.
Write its abelianization as $\mathbb Z^r\oplus A_1$.
Then
\[
 b_2=
 \begin{cases}
 d-1+\beta+2r,&L=0,\\
 d-1+\beta+L+r,&L>0,
 \end{cases}
 \qquad
 b_3=
 \begin{cases}
 2\beta+2r,&L=0,\\
 2\beta+2L,&L>0.
 \end{cases}                                               \tag{D.WP8}
\]
The complete integral groups are
\[
\begin{array}{c|c}
j&H_j(X;\mathbb Z)\\\hline
0,6&\mathbb Z\\
1&\mathbb Z^r\oplus A_1\\
2&\mathbb Z^{b_2}\oplus T\\
3&\mathbb Z^{b_3}\oplus T\\
4&\mathbb Z^{b_2}\oplus A_1\\
5&\mathbb Z^r .
\end{array}                                                \tag{D.WP9}
\]
In particular $\chi(X)=2d$. The executable
\path{code/manifolds_weighted_profiles.py} implements
all steps, including actual freeness, the full group
presentation and both Smith calculations.

\begin{proof}
\emph{The geometric meaning of the new weights.}
Iteration of the original lifted translation replaces
its line multiplier by the product of its first $d$
translates. If its initial zero is one component of
$I_n$, the new zero is its orbit sum, with repetitions.
It need not be a whole fibre. This is the prescribed
norm linearization, so no unsupported line-bundle
modification is being introduced.
Outside the finite replacements the new space is
the connected gamma cover defined by $S_d$ of the
old original or Mumford model. It follows directly
that the good-reduction monodromies are (D.WP2);
the same construction on their root discs permits
all the new free finite quotients.

At the unique rank-two zero the old valuation
lattice was the full short lattice. Multiplying
one long period by $d$ gives index $d$. The same
smooth triangular fan therefore has $d$ copies
of the old normalized $dP_6$ in its quotient.
At a zero-free rank-one site, write the old
monodromy as $I+\alpha\beta^*$ with $\alpha$
primitive and $\beta^*$ primitive. The cover
replaces $\beta^*$ by $\beta^*S_d$.
The exact factorization certificate retains all $436$
zero-free rank-one matrices as $\alpha\beta^*$, with
both vectors primitive. If $P$ is narrow, the content
of the first three entries of $\beta^*$ is exactly
$n_i$, and its gamma entry is prime to $n_i$. Thus
the new content is $\gcd(d,n_i)$. At every non-narrow
$P$ site that vertical content is exactly one, so
its new content is one for every $d$. These actual
vectors and entries are recorded in
\path{records/manifolds_weighted_rankone_certificates.json}.
Thus the respective component counts are precisely
those in the statement. This also follows directly
by lifting the component cycle through the gamma
cover. Zero-degree original multiplicative caps
have the same component count and the same global
kernel image; no relative diffeomorphism is asserted.

An original zero-degree additive cap is
$S^1_\delta$ times the mapping torus of the Kodaira
tree under $Q^d$. It has $\rho_i(d)$ fundamental
component cycles. At a nonzero-Chern original, $Q$
was already narrow and remains narrow; its
primitive Chern vector is unchanged. Its component
count and its principal-torus description from
D.2.32 or D.2.33 are therefore unchanged.
All these caps have torsion-free top homology and
Euler number zero. The rank-two cusp has Euler
number $2d$ by its normalization complex.

\emph{Why the regular integral lattice is free in the stated sector.}
For each original row, D.2.23 gives
\[
 \Lambda_{\{A_i\}}=\mathbb Z\gamma,\qquad
 (\Lambda^2\Lambda)_{\{A_i\}}=\mathbb Zs,
\]
with primitive functional (D.WP3) at $d=1$.
There is a finite exact certificate for all other
weights. Set $d=0$ in the polynomial matrices
(D.WP2), meaning that their first three gamma
entries are zero. For every one of the $172$ rows
the common degree-one relation matrix still has
rank three, the degree-two matrix has rank five,
and every prime in their nonunit Smith factors
divides $a$. The complete constant matrices and
Smith factors are retained in
\path{records/manifolds_weighted_zero_certificates.json}.

Here is why this finite certificate proves the
unbounded claim. At a prime not dividing $d$,
$S_d$ is an integral local isomorphism, so the
relations have the original saturated ranks.
At a prime dividing $d$, condition (D.WP1) makes
$a$ a unit, and reduction modulo that prime is
the just-certified $d=0$ matrix, still of full
rank three or five. Hence no prime divides a
nonunit Smith factor at the new weight.
The two common coinvariants are therefore free
of ranks one and one for every allowed $d$.
Moreover (D.WP3) is primitive since
$\gcd(a,db)=1$, so it is their actual integral
degree-two functional.

The graph Serre sequence consequently gives
\[
 H_1(V)=\mathbb Z\gamma\oplus\mathbb Z^{N-1},
 \qquad H_2(V)=\mathbb Zs\oplus K.                         \tag{D.WP10}
\]
The splitting used in (D.WP10) is marked. On the
original regular piece the integral fibre class
has value zero on every fixed-loop peripheral
sweep, by the unrolled comparison and the single
diagonal relation of D.2.23--26. Pull it back
through the actual gamma cover. It restricts to
$\omega_d$ and still has value zero on every
new fixed sweep: that sweep maps to the old
sweep of $S_d\lambda$.
Primitivity of $\omega_d$ makes this a genuine
integral splitting. This is a pullback of an
already marked class, not an assumption that
an invariant differential form descends integrally.

\emph{All local boundary-kernel images.}
For a finite-good matrix, all nonidentity linear
powers have the same fixed real two-plane.
Its projected lattice is
$P\Lambda=(N_h/h)\Lambda=F(H/h)\mathbb Z^2$.
The translation $v/m$ on that quotient has exact
order $m/t$. A nonidentity linear power can have
a fixed point exactly when its exponent is
divisible by $m/t$. Requiring all such exponents
to have identity linear part is $h\mid m/t$.
For the identity-linear powers one must additionally
have the full translation of order $m$, which is
the primitive-vector condition in (D.WP5).
This proves the stated freeness criterion for
arbitrary coordinates and orders.

The full normal Thom image is generated by the
four covering-torus loops and the affine meridian.
The cut of their swept normal tori into the regular
fibre is respectively $N_m\lambda$ and $v$.
Integration of the actual pulled-back marked
class from (D.WP10) gives respectively
$-\omega_d(\lambda,v)$ and zero, exactly as in
D.2.29. These are complete integral coordinates
in the free target (D.WP10), so a local torsion
ambiguity cannot be lost by this computation.
In particular no claim of integral form descent
on the finite quotient is needed.
This proves (D.WP6), including finite cap and
normal-line torsion.

For a Mumford or zero-degree original cap, every
new fixed normal torus bounds. Its image under
the covering is an old fixed normal torus; the
old bounding section-disc sweep lifts because
its loop is in the covering lattice. Equivalently
use the same toric or Neron-logarithm section
disc and its actual invariant circle actions.
These are all the normal projections, since
the Wang projection takes values in the full
fixed lattice.
After subtracting them, any further kernel
element comes from a fibre class. Rationally
its image is zero exactly on the monodromy-image
subspace. At the rank-two zero, the two primitive
angular normal tori exhaust the boundary kernel:
$H_2(\partial N)=\mathbb Z^6$ and
$H^3(N)=\mathbb Z^2$ by D.2.34. The regular-fibre
coinvariant of rank four therefore injects. At a
rank-one cap the nodal cycle times $T^2$ supplies
four independent fibre images (whole fibre, two
mixed surviving-circle tori, and transverse area).
At an original zero-degree additive cap the
Kodaira-tree permutation model supplies the whole
fibre and the transverse area, the full two
dimensional coinvariant. These explicit models
prove the required ranks for every weight.
The functional $\omega_d$ annihilates that
subspace. Since the global fibre coinvariant
in (D.WP10) is free, every such fibre-kernel
class has zero global image, even when it is
torsion in the boundary. This proves the entire
global image (D.WP7).

For a nonzero-Chern original, the marked logarithm
vectors scale by $d$. The exact area-kernel
calculations of D.2.32--33 therefore multiply
their fibre values by $d$, giving $dh_i\mathbb Zs$.
Their complete fixed normal projections still
bound, and the third multiplicative correction
is still annihilated by $\omega_d$.
This proves the remaining column, integrally.

\emph{The relative sequence and the free component contribution.}
First attach only the finite and regular logarithmic
caps to $V$, and call the result $E$. There is at
least one unfilled cap, the rank-two zero.
The relative degree-two group of each attached
disc bundle is $\mathbb Z$, whose boundary has
its nonzero multiplicity in its independent base
meridian coordinate. These relative maps into
$H_1(V)$ are injective. Thus $H_2(E)$ is exactly
the quotient of (D.WP10) by their full Thom
images, including when the cap has torsion.

Now attach all the remaining original and Mumford
caps $N_j$. Their relative degree-two groups are
\[
 H_2(N_j,\partial N_j)=H^4(N_j),
\]
free on their fundamental component cycles.
The degree-three boundary images are exactly
the kernels already computed above. Consequently
the pair sequence has the form
\[
 0\longrightarrow\operatorname{coker}B
 \longrightarrow H_2(X)
 \longrightarrow
 \ker\!\left\{\bigoplus_j H^4(N_j)\to H_1(E)\right\}
 \longrightarrow0.                                      \tag{D.WP11}
\]
Its last group is a subgroup of a free group,
so the sequence splits as a sequence of abelian
groups. This avoids all potentially nonprimitive
degree-one boundary splittings.

If there are $s$ unfilled caps, $H_1(E)$ has free
rank $s$: the finite/regular power relations
remove their independent base directions.
The final rank is $r$, so the map in (D.WP11)
has rank $s-r$. Its domain has rank equal to
the sum of the component counts. Subtracting
one per cap gives
\[
 \operatorname{rank}\ker
   =d-1+\beta+L+r.                                        \tag{D.WP12}
\]
This is an exact rank calculation in a free
relative group; torsion meridians can change
its index but cannot change its freeness or
introduce an extension in (D.WP11).

Over the rationals (D.WP2) is conjugate to the
original monodromy. Thus the fixed-sweep graph
has the same single diagonal relation as in
D.2.23--26. With no nonzero-Chern original,
the remaining scalar equation is the actual
Seifert Euler equation, so
$\operatorname{corank}_{\mathbb Q}B=r$.
If $L>0$, a nonzero fibre multiple is killed,
and then every fixed sweep is killed rationally;
the corank is zero. Combining this with
(D.WP12) proves $b_2$ and its full torsion $T$.

The fundamental-group argument is also unchanged
integrally. The primitive angular vanishing
vectors are vertical, so they survive in the
cover without a multiple. Restrictions of
$A_i(d)$ to the vertical lattice are the
original restrictions. Their normal closures
therefore kill the same entire vertical
lattice as in D.2.23, including its original
star thimble when required. Each unchanged
cap has an actual transverse origin disc, so
adds no gamma twist beyond the recorded
peripheral clutch. The finite/regular power
relations then give exactly the full Seifert
group, not merely its abelianization.
Finally $\chi=2d$ and integral duality give
the remainder of (D.WP8)--(D.WP9).
\end{proof}

The program verifies all $172$ weight-zero
prime-support certificates and computes $1626$
full weighted graph presentations on the $495$
patterns, of which $1131$ have weight greater
than one. The all-parameter proof is the
prime-local argument and marked relative
calculation above, not a finite weight search.
Weights sharing a factor with $a$, additional
star logarithmic Mumford choices and arbitrary
initial pairs not obtained by these translation
iterates remain outside this theorem. Its
homology algorithm makes no new smooth
classification assertion.
